\section{Curvature at the radius and cap scales}\label{sec:cap}

\begin{proposition}\label{prop:full-curvature}
For every flow with the choices made above, and \(q\) sufficiently large,
there is a constant \(C\), allowed to depend on that flow, such that
\begin{equation}\label{eq:full-curvature}
(r^2+\theta^2)|\Rm|\le C
\end{equation}
on \(M\times[0,T)\).
\end{proposition}

We prove the two scale bounds separately. We use Shi's curvature derivative
estimates~\cite{Shi1989} in their global and local forms; precise statements
and proofs are also given in~\cite[Theorems~1.4.1--1.4.2]{CaoZhu2006}.
A uniform curvature
bound on a parabolic neighborhood gives a bound for \(|\nabla\Rm|\) on a
smaller neighborhood away from its initial time. The constants depend on
the dimension, curvature bound, spatial margin, and elapsed time, and do
not require an injectivity-radius lower bound. In each application below,
the requisite parabolic neighborhood is established before the derivative
estimate is used.

Stolarski proves a radius curvature bound in the inner region
\cite[Proposition~5.1(3)]{Stolarski2019}. We first establish a global
radius bound and then the cap-scale bound in
Proposition~\ref{prop:full-curvature}.
If the radius bound failed, both rescaled sphere radii
would diverge while their arclength derivatives remained bounded. A
persistent nonzero radial sectional curvature would then force a large
change in one of those derivatives. For the cap bound, we first need a
lower bound for the radius of the pole orbit and a quantitative estimate
showing that \(r_s\) approaches one near the pole. The same radial
integration argument can then be used at the cap scale.

\begin{lemma}\label{lem:r-curvature}
There is a time-independent constant \(C\) such that \(r^2|\Rm|\le C\).
\end{lemma}

\begin{proof}
Suppose otherwise. Choose record points \((x_i,t_i)\), with \(t_i\uparrow T\),
and put \(r_i=r(x_i,t_i)>0\), \(Q_i=|\Rm|(x_i,t_i)\), so that
\begin{equation}\label{eq:r-record}
r_i^2Q_i=\max_{M\times[0,t_i]}r^2|\Rm|\longrightarrow\infty.
\end{equation}
Such points exist by compactness on each closed time slab. Boundedness of
\(r\) implies \(Q_i\to\infty\). By the \(\Sp^q\) component of the
metric evolution,
\[
|\dt|_x r^2|=2|\lambda_q|r^2\le C_mr^2|\Rm|.
\]
The identity extends to the poles since \(r^2\) is smooth. For fixed
sufficiently small \(b>0\), Equation~\eqref{eq:r-record} gives
\begin{equation}\label{eq:r-time-variation}
|r(x,t)^2-r(x,t_i)^2|\le r_i^2/16
\quad\text{for }t_i-b/Q_i\le t\le t_i,
\end{equation}
at every fixed manifold point \(x\). These time intervals lie in the
flow for all large \(i\).

Rescale by \(Q_i\) and translate \(t_i\) to zero. On a fixed ball of
radius \(R\) about \(x_i\) in the rescaled metric at time \(-b\), the
Lipschitz bound for \(r\), Equation~\eqref{eq:r-time-variation}, and
\(r_i\sqrt{Q_i}\to\infty\) imply \(r\ge r_i/2\) throughout
\([-b,0]\), for large \(i\). The record bound thus gives
\(|\Rm|/Q_i\le4\) on this fixed spatial domain throughout the interval.
Integrating the metric evolution makes the metrics uniformly comparable
there, with a length-comparison factor \(E\) independent of \(R\) and \(i\).
Choose \(R>E+2\). A final-time curve of length at most one cannot leave
this initial ball: up to its first exit its initial length would be at
least \(R\), and hence its final length at least \(R/E>1\).
The final unit ball therefore lies inside the initial ball, with a fixed
positive margin measured in the initial metric. Take \(b\) below the
dimension-dependent time threshold in the local derivative estimate.
That estimate gives a uniform
bound for the final rescaled \(|\nabla\Rm|\) on that ball.

In the final rescaled metric the warping functions are
\(F_1=\sqrt{Q_i}\phi\) and \(F_2=\sqrt{Q_i}r\). Both tend to infinity
at the center by Lemma~\ref{lem:annuli}. Their derivatives with respect
to rescaled arclength are the original \(\phi_s,r_s\), hence uniformly
bounded. Therefore both functions tend uniformly to infinity on any fixed
short radial interval about the center. Such intervals exist: the
Lipschitz bound makes the rescaled distance to either pole tend to
infinity. Crossing the equator presents no obstruction.

The rescaled tangential curvatures \(j,\ell,\mu\) tend to zero. Since
the rescaled norm of curvature is one at the center, the identity
\begin{equation}\label{eq:curvature-norm}
|\Rm|^2=4\left(ph^2+qh_q^2+\frac{p(p-1)}2j^2
 +\frac{q(q-1)}2\ell^2+pq\mu^2\right)
\end{equation}
shows that a radial curvature \(-F_a''/F_a\) has absolute value bounded
below by a positive dimension-dependent constant. Along a radial geodesic,
the radial vector and a fixed sphere-tangent vector divided by its warping
factor are parallel. The derivative estimate therefore keeps that same
sectional curvature of one sign and bounded away from zero on a fixed
shorter interval. On that interval \(F_a\to\infty\) uniformly. Integrating
\(F_a''\) forces an unbounded change in \(F_a'\), contradicting the
slope bound.
\end{proof}

The next proof quantifies the scalar-sign obstruction in
\cite[Proposition~5.3]{Stolarski2019}. There, the strict cone gap and
monotone logarithmic slope rule out an ancient limit with nonnegative
scalar curvature. Here we use the lower bound \(R\ge-C\) on the given
flow and integrate over a finite annulus whose logarithmic width would
diverge if the pole orbit were too small.

\begin{lemma}[Size of the pole orbit]\label{lem:pole-size}
At all sufficiently late times,
\begin{equation}\label{eq:pole-size}
c\theta\le\phi(0,t)\le C\theta.
\end{equation}
Moreover \(\phi(r,t)\ge c\theta\) for \(0\le r\le Y_1\theta\).
\end{lemma}

\begin{proof}
The upper bound follows from Lemma~\ref{lem:gap} and \(f=ru_r\ge0\).
For the lower bound, the scalar curvature computed from
\eqref{eq:sectional} is
\begin{equation}\label{eq:scalar-f}
\begin{split}
r^2R={}&p(p-1)(r/\phi)^2+q(q-1)(1-z)\\
&-pz\bigl[(p+1)f^2+(2q-2)f+2rf_r\bigr]
 -(pf+q)rz_r.
\end{split}
\end{equation}
For example,
\(\phi_{ss}/\phi=z(u_{rr}+u_r^2)+z_ru_r/2\) and
\(r_{ss}=z_r/2\); substituting \(u_r=f/r\) gives this formula directly.

If the lower bound failed, choose times tending to \(T\) for which
\(L_0=\phi(0,t)=o(\theta)\). The inequalities \(r_s\ge B\) and
\(|\phi_s|\le C\) imply
\[
\phi(L_0,t)/L_0\le1+C/B.
\]
Let \(r_*=(L_0\theta)^{1/2}\). Since
\(\partial_{\log r}\log(\phi/r)=f-1\), the lower ratio bound and the
last upper bound give
\[
0\le\int_{L_0}^{r_*}(1-f)\,\frac{dr}{r}\le C.
\]
The integrand is nonnegative and nonincreasing, because \(0\le f\le1\)
and \(f_r\ge0\). Thus
\[
0\le1-f(r_*,t)\le\frac{C}{\log(r_*/L_0)}\longrightarrow0.
\]
Consequently \(f=1+o(1)\) uniformly on \([r_*,\theta]\).

At \(f=1\), the nondifferentiated terms in the first two lines of
\eqref{eq:scalar-f}, before the \(z_r\) term, reduce to
\[
\frac{p(p-1)}{(\phi/r)^2}+q(q-1)-m(m-1)z.
\]
They vanish at \(\phi/r=A/B\), \(z=B^2\). By the fixed strict gap
\eqref{eq:strict-gap} and \(z\ge B^2\), they are at most
\(-\kappa\) for some fixed \(\kappa>0\). The term involving \(f_r\)
is nonpositive. The uniform error \(f-1=o(1)\) is harmless with the flow
and dimension fixed. Using \(R\ge-C\) from Lemma~\ref{lem:gradients},
we conclude, after reducing \(\kappa\) if necessary, that
\[
(pf+q)rz_r\le-\kappa+C\theta^2\le-\kappa/2
\quad(r_*\le r\le\theta).
\]
Since \(pf+q\le m\), this implies \(rz_r\le-\kappa/(2m)\).
Integration contradicts \(B^2\le z\le1\), because
\(\log(\theta/r_*)=\tfrac12\log(\theta/L_0)\to\infty\).
This proves the pole lower bound. Monotonicity of \(\phi\) extends it
to the entire indicated inner interval.
\end{proof}

\begin{lemma}[Radial slope at the cap scale]\label{lem:cap-slope}
For sufficiently large \(q\), at late times,
\begin{equation}\label{eq:cap-slope}
0\le1-v\le C r/\theta\qquad(0\le r\le\theta).
\end{equation}
\end{lemma}

\begin{proof}
Let \(w=1-v\). Equation~\eqref{eq:sideways} gives
\begin{equation}\label{eq:w-operator}
\left(\partial_t-z\partial_{rr}-\frac{q-1-z}{r}\partial_r
       +\frac{(q-1)v(1+v)}{r^2}\right)w
=pv^3u_r^2\le C\theta^{-2}.
\end{equation}
The last inequality uses \(u_r=\phi_s/(\phi v)\),
Lemma~\ref{lem:pole-size}, and \(B\le v\le1\).
Applying this operator to \(r/\theta\) gives exactly
\[
\frac{\sigma r}{\delta\theta}
+\frac{(q-1)[v(1+v)-1]+z}{r\theta}.
\]
For large \(q\), \(B(1+B)>1\). Thus this expression is bounded below
by a fixed positive multiple of \(\theta^{-2}\) when \(r\le\theta\).
A sufficiently large multiple of \(r/\theta\) dominates the source,
the data at \(r=\theta\), and the data on a fixed late starting slice.

To justify comparison at the pole, fix a compact time slab. Smooth polar
expansions give \(w=O(r^2)\) uniformly on that slab. The proposed upper
barrier therefore dominates on a sufficiently small inner boundary
\(r=\eta\). Compare on \(\eta\le r\le\theta(t)\) and let
\(\eta\downarrow0\). The auxiliary radius may depend on the slab, but
the barrier multiplier does not. This proves~\eqref{eq:cap-slope}.
\end{proof}

\begin{proof}[Proof of Proposition~\ref{prop:full-curvature}]
It remains to bound \(\theta^2|\Rm|\). If this were unbounded, choose
record points with \(Q_i=|\Rm|(x_i,t_i)\) such that
\[
\theta_i^2Q_i=\max_{M\times[0,t_i]}\theta(t)^2|\Rm|(x,t)\to\infty,
\qquad \theta_i=\theta(t_i).
\]
Since \(\theta\) decreases, for every \(t\le t_i\) we have
\[
|\Rm|(x,t)\le\frac{\theta_i^2}{\theta(t)^2}Q_i\le Q_i.
\]
The global derivative estimate therefore bounds \(|\nabla\Rm|\) at
the final time after rescaling by \(Q_i\), uniformly in \(i\).
For all large \(i\), a fixed backward interval in rescaled time lies
inside the flow.

Lemma~\ref{lem:r-curvature} gives \(r_i\sqrt{Q_i}\le C\), whereas
\(\theta_i\sqrt{Q_i}\to\infty\). Any fixed short outward radial
interval from the center consequently lies in \(r\le\theta_i\).
On such intervals the rescaled first warping function satisfies
\(\sqrt{Q_i}\phi\ge c\sqrt{Q_i}\theta_i\to\infty\).
The rescaled second warping function is bounded above on each such
interval, and its arclength derivative satisfies
\[
v\longrightarrow1
\]
uniformly by Lemma~\ref{lem:cap-slope}.
If the center is on a pole orbit, choose any outward radial ray.

Move a fixed sufficiently small positive distance along the ray. The
derivative estimate keeps the curvature norm at least \(1/2\), while
the rescaled radius is now bounded below by a fixed positive number.
At this new center the three tangential curvatures tend to zero. By
\eqref{eq:curvature-norm}, some radial curvature has nonzero magnitude
and constant sign on a fixed subsequent interval. For
\(F=\sqrt{Q_i}\phi\), integration of \(F''\) contradicts bounded
\(F'\) because \(F\to\infty\). For \(F=\sqrt{Q_i}r\), it contradicts
the uniform convergence \(F'=v\to1\), because \(F\) stays bounded below
by a positive constant on that interval. Both alternatives are impossible.
Together with Lemma~\ref{lem:r-curvature}, this proves~\eqref{eq:full-curvature}.
\end{proof}
